\documentclass[10pt,a4paper]{amsart}
\usepackage[margin=19mm]{geometry}
\usepackage{amsmath,amssymb,mathtools}
\usepackage[scr=rsfs,cal=euler]{mathalfa}
\usepackage{xfrac}
\usepackage[unicode,hidelinks,bookmarks=false]{hyperref}
\newcommand{\ie}{i.e.\ }
\newcommand{\cf}{cf.\ }
\newcommand{\resp}{respectively\ }
\newcommand{\Subsection}{\subsection}
\providecommand{\nobreakdash}{\nobreak\hbox{-}}
\setlength{\emergencystretch}{2em}
\multlinegap0pt
\usepackage{mathtools}
\usepackage{bm}
\usepackage{amssymb}
\usepackage{dsfont}
\usepackage{xcolor}
\usepackage{tikz}
\usepackage[shortlabels]{enumitem}
\newtheorem{lemma}{Lemma}
\usepackage{cleveref}
\crefname{lemma}{Lemma}{Lemmas}
\newcommand{\R}{\mathbb{R}} % reelle
\DeclareMathOperator{\ch}{ch}
\DeclareMathOperator{\dcs}{sp}
\DeclareMathOperator{\dct}{tm}
\newcommand*\diff{\mathop{}\!\mathrm{d}}
\newcommand{\fpart}{\Pi} % \Pi, P, \eta
\newcommand{\fsp}{f_\mathrm{sp}}
\newcommand{\interp}{\overline{x}} % for lem:NFaux
\DeclareMathOperator{\lf}{lf}
\DeclareMathOperator{\nd}{nd}
\newcommand{\rfy}[1]{\widetilde{{#1}}} % for an Rd-version of something defined for the torus.
\DeclareMathOperator{\rt}{rt}
\title{The Brownian Spatial Coalescent}
\author{Peter Koepernik}
\date{Source: arXiv:2401.08557v4 (CC BY 4.0)}
\begin{document}
\maketitle

\noindent\emph{Source and licence.} The Brownian Spatial Coalescent. Version arXiv:2401.08557v4 (CC BY 4.0), distributed by arXiv under the Creative Commons Attribution 4.0 International licence, http://creativecommons.org/licenses/by/4.0/, as recorded in the arXiv licence metadata for this exact version and shown on the abstract page https://arxiv.org/abs/2401.08557v4. Full licence notice: \texttt{licenses/arxiv-2401.08557-CC-BY-4.0.txt}.\par\smallskip

\noindent\textit{Editorial scope.} Counted unit: the article's lemma lem:NFaux (source0.tex lines 3866--3883). The article's complete original proof of it is delivered with it (source0.tex lines 3884--3927, one \texttt{proof} environment of 44 lines). No mathematics is written, repaired or re-derived: the statement and the proof are the article's own lines, in the article's own notation, and no retained line is substituted or re-flowed.  Retained uncounted as the source-local context the proof needs: lines 3841--3845.  Nothing was omitted from any retained span, so the package carries no omission record.  No retained line contains a citation, so the wrapper carries no bibliography and no bibliography entry.  No retained line contains a figure environment, an image-inclusion command or drawing code, so no image is delivered and the package declares no asset dependency.  Editorial formatting only, in addition: an AMS \texttt{amsart} preamble that restates the article's own declarations and macros used by the retained lines, namely the environments \texttt{lemma} on separate counters, together with the standard packages those lines need.  The retained environments print in this document as Lemma 1 in order of appearance, whereas the article prints the same items with the numbering of its own full document.  The complete native source of the article as distributed in the arXiv e-print for this exact version is bundled unchanged as \texttt{sources/153037/source0.tex}.
\begin{lemma}\label{lem:heatkerneltrick}
    Let $x_1, \ldots ,x_n \in \R^d$ and $s_1, \ldots ,s_n > 0$. Then, \[
        \prod_{i=1}^n \rfy{p}_{s_i}(x_i-z) = \frac{\rfy{p}_s(z-\overline{x})}{\rfy{p}_s(0)} \prod_{i=1}^n \rfy{p}_{s_i}(x_i - \interp),
    \] where $\frac{1}{s} = \sum_{i=1}^n \frac{1}{s_i}$ and $\interp = \sum_{i=1}^n \frac{s}{s_i}x_i$.
\end{lemma}

\begin{lemma}\label{lem:NFaux}
    For a non-trivial forest $F\in \mathbb{F}$ and $\boldsymbol x \in \rfy{E}^{\lf(F)}_\circ $, $\tau \in \dct(F)$,
    \begin{align}\label{eq:NFaux}
        \int_{\rfy{\dcs}(F)} \rfy{\fsp}(\xi \,\vert\,\tau,\boldsymbol x) \diff \xi = \prod_{v\in \nd^\circ (F)} \frac{1}{\rfy{p}_{r_v}(0)} \prod_{u\in \ch_F(v)} \rfy{p}_{r_u + \tau_v - \tau_u} (\interp_v -\interp_u),
    \end{align}
    where $\interp_u = \boldsymbol x_u$ and $r_u = 0$ for $u\in \lf(F)$ and, inductively for $v\in \nd^\circ (F)$, \[
        \frac{1}{r_v} \coloneqq \sum_{u\in \ch_F(v)} \frac{1}{r_u+\tau_v - \tau_u},\qquad \interp_v \coloneqq \sum_{u\in \ch_F(v)} \frac{r_v}{r_u+\tau_v - \tau_u} \interp_u.
    \] The following hold.
    \begin{enumerate}
        \item If $v\in \nd^\circ (F)$ with $\ch(v) \subset \lf(F)$, then \[
                r_v = \frac{\tau_v}{|\ch_F(v)|},\qquad \interp_v = \frac{1}{|\ch_F(v)|} \sum_{u\in \ch_F(v)} \boldsymbol x_u.
        \]
        \item There is $c = c(F) > 0$ such that $c \tau_v \le r_v < \tau_v$ for all $v\in \nd^\circ (F)$.
        \item If $v\in \nd^\circ (F)$ then \[
                \max_{u\in \ch(v)} |\interp_u-\interp_v| \ge \frac{1}{2}\max_{u,u'\in \ch(v)}|\interp_u-\interp_{u'}|.
        \]
    \end{enumerate}
\end{lemma}

\begin{proof}
    \Cref{eq:NFaux} follows by applying \cref{lem:heatkerneltrick} inductively from leaves to root, which is tedious but straightforward.
    %\emph{Full proof commented out.}
    %Throughout the proof we abbreviate $\ch \coloneqq \ch_F$ if $F$ is clear from context. We first prove that
    %\begin{align*}
    %    \int \fsp(\xi \,\vert\,\tau,\boldsymbol x) \smashoperator[r]{\prod_{u\in \nd(F) \setminus \rt(F)}}\, \diff \xi_u
    %    &= \prod_{v\in \nd^\circ (F)\setminus \rt(F)} \frac{1}{\rfy{p}_{r_v}(0)} \prod_{u\in \ch(v)} \rfy{p}_{r_u+\tau_v-\tau_u}(\interp_u-\interp_v)\\
    %    &\phantom{=} \times \prod_{v\in \nd^\circ (F) \cap \rt(F)} \frac{\rfy{p}_{r_v}(\xi_v-\interp_v)}{\rfy{p}_{r_v}(0)} \prod_{u\in \ch(v)} \rfy{p}_{r_u+\tau_v-\tau_u}(\interp_u-\interp_v).
    %\end{align*}
    %Then integrating out $\xi_v,\,v\in \rt(F)$ gives the claim. We do induction over $|F|$. Suppose $F = (\fpart_0,\fpart_1)$, then
    %\begin{align*}
    %    \fsp(\xi \,\vert\,\tau,\boldsymbol x)
    %    &= \prod_{v \in \fpart_1\setminus \fpart_0} \prod_{u\in \ch(v)} \rfy{p}_{\tau_v}(\interp_u - \xi_v)\\
    %    &= \prod_{v\in \rt(F)} \frac{\rfy{p}_{r_v}(\xi_v-\interp_v)}{\rfy{p}_{r_v}(0)} \prod_{u\in \ch(v)} \rfy{p}_{\tau_v}(\interp_u - \interp_v).
    %\end{align*}
    %Now suppose $F = (\fpart_0, \ldots ,\fpart_m)$ and we proved the claim for $|F| \le m$. Denote $F' \coloneqq (\fpart_0, \ldots ,\fpart_{m-1})$,
    %\begin{align*}
    %    \fsp(\xi \,\vert\,\tau,\boldsymbol x)
    %    &= \fsp^{F'}(\xi \,\vert\,\tau,\boldsymbol x) \prod_{v\in \fpart_m\setminus \fpart_{m-1}} \prod_{u\in \ch(v)} \rfy{p}_{\tau_v-\tau_u}(\xi_v-(\boldsymbol x \xi)_u)\\
    %    &=
    %\end{align*}
    %Integrating over $\xi_u$ for $u\in \nd(F') \setminus \rt(F') = \fpart_0\cup\ldots \cup \fpart_{m-1}$ gives, by induction hypothesis,
    %\begin{align*}
    %    &\prod_{v\in \nd^\circ (F')\setminus \rt(F')} \frac{1}{\rfy{p}_{r_v}(0)} \prod_{u\in \ch(v)} \rfy{p}_{r_u+\tau_v-\tau_u}(\interp_u-\interp_v)\\
    %    & \times \prod_{v\in \fpart_{m-1}} \frac{\rfy{p}_{r_v}(\xi_v-\interp_v)}{\rfy{p}_{r_v}(0)} \prod_{u\in \ch(v)} \rfy{p}_{r_u+\tau_v-\tau_u}(\interp_u-\interp_v)\\
    %    & \times \prod_{v\in \fpart_m\setminus \fpart_{m-1}} \prod_{u\in \ch(v)} \rfy{p}_{\tau_v-\tau_u}(\xi_v-(\boldsymbol x \xi)_u).
    %\end{align*}
    %etc.

    \begin{enumerate}
        \item is obvious from the definitions.
        \item If $v$ is not a leaf and $r_u \le \tau_u$ holds for all children, then
            \begin{align*}
                \frac{1}{r_v} \ge \sum_{u\in \ch_F(v)} \frac{1}{\tau_u + \tau_v - \tau_u} = \frac{|\ch_F(v)|}{\tau_v} \ge \frac{2}{\tau_v} > \frac{1}{\tau_v}.
            \end{align*}
            If $v$ is not a leaf and $r_u \ge c \tau_u$ holds for all children, then
            \begin{align*}
                \frac{1}{r_v} \le \sum_{u\in \ch_F(v)} \frac{1}{c \tau_u + \tau_v - \tau_u} \le \sum_{u\in \ch_F(v)} \frac{1}{c(\tau_u + \tau_v - \tau_u)} = \frac{|\ch_F(v)|}{c} \frac{1}{\tau_v}.
            \end{align*}
        \item If $v\in \nd^\circ (F)$ and $u_1,u_2\in \ch(v)$, then \[
                |\interp_{u_1}-\interp_{u_2}| \le |\interp_{u_1} - \interp_v| + |\interp_v - \interp_{u_2}| \le 2 \max_{u\in \ch(v)}|\interp_u-\interp_v|.
        \]
    \end{enumerate}
\end{proof}

\end{document}
